
\section{Experiments}
To validate that our system—comprising the fine-tuned model and its integrated CEO—can effectively enhance collaboration and reasoning in MAS, we conduct experiments using both state-of-the-art (SOTA) open-source and closed-source LLMs on AgentVerse across tasks in general understanding, mathematical reasoning, and coding. Additional investigations are conducted to investigate the emerging behavior and scaling performance of our method.

\subsection{Experimental Details}\label{sec:experimental_details}
\paragraph{LLMs.} We evaluate both reasoning-oriented and non-reasoning LLMs to fully understand the effect of collaboration and reasoning in MAS. The primary baselines include Qwen2.5-32B-Instruct (abbreviated as \textbf{Qwen2.5}) \citep{qwen25} and s1.1-32B \citep{s1}. Both M1-32B and s1.1-32B are fine-tuned from Qwen2.5; s1.1-32B additionally utilizes questions from Simple-Scaling \citep{s1} using DeepSeek-R1 reasoning traces in a single-agent setting. We also include DeepSeek-V3 \citep{deepseek_v3} and DeepSeek-R1 \citep{r1} as strong open-source baselines. For closed-source models, we use o3-mini (medium) \citep{o3-mini} and GPT-4o (GPT-4o-2024-08-06) \citep{gpt4o}.

\paragraph{Tasks.} To conduct a comprehensive evaluation, we focus on three critical domains: \ding{182} \textbf{General Understanding}: We use GPQA-Diamond (abbreviated as \textbf{GPQA}) \citep{gpqa}  to evaluate the general knowledge and  Commongen-Challenge (abbreviated as \textbf{Commongen}) \citep{self-refine} to evaluate sentence writing and response readability.  GPQA-Diamond contains 198 PhD-level science questions from Biology, Chemistry, and Physics. We report the percentage of questions answered correctly (zero-shot). In Commongen-Challenge, the agent is required to generate a coherent and grammatically correct paragraph using as many of the 20 given concepts as possible. The benchmark consists of 200 concept lists, and we report the average percentage of covered concepts. \ding{183} \textbf{Mathematical Reasoning}: We evaluate on two widely used challenging math benchmarks: AIME2024 \citep{aime2024} and MATH-500 \citep{math}. AIME2024 includes 30 problems from the 2024 American Invitational Mathematics Examination (AIME), while MATH-500 is a curated benchmark of competition-level math problems with varying difficulty. The zero-shot accuracy, i.e., the percentage of correctly solved problems, is reported. \ding{184} \textbf{Coding}: We evaluate code generation ability using HumanEval \citep{humaneval} and MBPP-Sanitized (abbreviated as \textbf{MBPP-S}) \citep{mbpp}. HumanEval consists of 164 Python programming problems designed to test the ability to generate functionally correct code from natural language specifications. MBPP-Sanitized contains 257 introductory Python programming problems that cover a broad range of algorithmic and functional challenges. For both benchmarks, we report the zero-shot Pass@1 accuracy.

\paragraph{Training and Evaluation.} We perform SFT on Qwen2.5 using the M500 dataset for 5 epochs with a learning rate of 1e-5, resulting in our model M1-32B. Training is conducted on 8 NVIDIA A100 GPUs using FlashAttention \citep{flashattention} and DeepSpeed \citep{deepspeed} within the LLaMA-Factory framework \citep{llamafactory}. Evaluation is carried out using the open-source MAS AgentVerse with a default total agent number of 5, critic iteration number of 3, and total iteration number of 2. The final response generated by the MAS is used for evaluation. All main results are averaged over three runs. The prompts used for all agents in the mathematical reasoning tasks are detailed in Appendix~\ref{sec:appendix_prompts} for reproducibility, with prompts for other tasks available in the accompanying code.

\subsection{Main Results}

\begin{table}[t]
\centering
\resizebox{1\textwidth}{!}{%
\begin{tabular}{lcccccc}
\toprule
\multirow{2}{*}{\textbf{Model}} & \multicolumn{2}{c}{\textbf{General Understanding}} & \multicolumn{2}{c}{\textbf{Mathematical Reasoning}} & \multicolumn{2}{c}{\textbf{Coding}} \\
\cmidrule(lr){2-3} \cmidrule(lr){4-5} \cmidrule(lr){6-7}
 & GPQA & Commongen & AIME2024 & MATH-500 & HumanEval & MBPP-S  \\
\midrule
\multicolumn{7}{c}{\textit{Non-Reasoning Models}} \\
\midrule
Qwen2.5 & 50.2 & 96.7 & 21.1 & 84.4 & 89.0 & 80.2 \\
DeepSeek-V3           & \textbf{58.6} & \textbf{98.6} & \textbf{33.3} & \textbf{88.6} & 89.6 & 83.9 \\
GPT-4o                & 49.2 & 97.8 & 7.8  & 81.3 & \textbf{90.9} & \textbf{85.4} \\
\midrule
\multicolumn{7}{c}{\textit{Reasoning Models}} \\
\midrule
s1.1-32B              & 58.3 & 94.1 & 53.3 & 90.6 & 82.3 & 77.4 \\
DeepSeek-R1           & \textbf{75.5} & 97.2 & 78.9 & \textbf{96.2} & \textbf{98.2} & 91.7 \\
o3-mini               & 71.3 & \textbf{99.1} & \textbf{84.4} & 95.3 & 97.0 & \textbf{93.6} \\
\midrule
M1-32B (Ours)               & 61.1 & 96.9 & 60.0 & 95.1 & 92.8 & 89.1 \\
M1-32B w. CEO (Ours)         & 62.1 & 97.4 & 62.2 & 95.8 & 93.9 & 90.5 \\
\bottomrule
\end{tabular}}
\caption{Performance comparison on general understanding, mathematical reasoning, and coding tasks using strong reasoning and non-reasoning models within the AgentVerse framework. Our method achieves substantial improvements over Qwen2.5 and s1.1-32B on all tasks, and attains performance comparable to o3-mini and DeepSeek-R1 on MATH-500 and MBPP-S, demonstrating its effectiveness in enhancing collaborative reasoning in MAS.
}
\label{tab:main_results}
\end{table}

The experimental results comparing our method and baseline models across general understanding, mathematical reasoning, and coding tasks are presented in Table \ref{tab:main_results}. Several key findings emerge from these results:
\begin{itemize}
    \item Our proposed method achieves substantial performance improvements across all evaluated tasks relative to Qwen2.5, demonstrating that the integration of M1-32B and the CEO agent effectively enhances general question answering, writing, mathematical reasoning, and coding capabilities within MAS. Specifically, M1-32B w. CEO improves performance by 12\%, 41\%, 11\%, and 10\% on GPQA, AIME2024, MATH-500, and MBPP-S, respectively, compared to Qwen2.5. Moreover, our method achieves comparable performance with SOTA open-source and closed-source models, such as DeepSeek-R1 and o3-mini, on MATH-500, Commongen, and MBPP-S, underscoring the effectiveness of our approach.

    \item Our approach significantly enhances collaborative reasoning in MAS compared to the Simple-Scaling \citep{s1}. For instance, M1-32B with CEO outperforms s1.1-32B by 4\% and 9\% on GPQA and AIME2024, respectively. Additionally, s1.1-32B experiences performance degradation in coding tasks compared to Qwen2.5, likely due to the limited coding examples in the Simple-Scaling dataset. In contrast, our method notably enhances coding performance, highlighting its advantage over Simple-Scaling. Both M1-32B and s1.1-32B are trained on samples derived from the Simple-Scaling dataset; thus, the observed improvements indicate that multi-agent collaborative reasoning traces are more effective than single-agent reasoning traces in enhancing LLM capabilities within MAS.

    \item The introduction of the CEO agent consistently improves the performance of M1-32B across all tasks, highlighting that collaborative reasoning in MAS is effectively scaled when guided by an M1-32B-based CEO agent.
\end{itemize}


\begin{figure}[t]
\centering
\begin{tcolorbox}[
  enhanced,
  title=Aha Moment in MAS,
  colback=white,
  colframe=black,
  colbacktitle=gray!20,
  coltitle=black,
  fonttitle=\bfseries,
  boxrule=0.8pt,
  sharp corners,
  left=2pt, right=2pt, top=1pt, bottom=1pt
]
\scriptsize

% Question
\begin{tcolorbox}[
  enhanced,
  colback=white,
  colframe=royalblue,
  boxrule=0.5pt,
  sharp corners,
  left=1pt, right=1pt, top=1pt, bottom=1pt
]
\textcolor{royalblue}{\textbf{Question}} \\
Let $ABC$ be a triangle inscribed in circle $\omega$. The tangents to $\omega$ at $B$ and $C$ intersect at point $D$, and $\overline{AD}$ intersects $\omega$ again at $P$. If $AB=5$, $BC=9$, and $AC=10$, $AP$ can be written as $\frac{m}{n}$ in lowest terms. Find $m + n$.
\end{tcolorbox}

% Problem Solver
\begin{tcolorbox}[
  enhanced,
  colback=white,
  colframe=olivegreen,
  boxrule=0.5pt,
  sharp corners,
  left=1pt, right=1pt, top=1pt, bottom=1pt
]
\textcolor{olivegreen}{\textbf{Problem Solver}} \\
... Place $B = (0,0)$, $C = (9,0)$, compute power of $D$: \(DB^2 = \left(\frac{225}{22}\right)^2\). Use the power of a point: \(DA \cdot DP = DB^2\), solve for \(DP\), then compute \(AP = DP - DA\). Final result: \(AP = \frac{22}{1} \Rightarrow m+n = \boxed{23}\).
\end{tcolorbox}

% CEO
\begin{tcolorbox}[
  enhanced,
  colback=white,
  colframe=orange,
  boxrule=0.5pt,
  sharp corners,
  left=1pt, right=1pt, top=1pt, bottom=1pt
]
\textcolor{orange}{\textbf{CEO}} \\
\textless think\textgreater ...\textcolor{red}{Wait, let's see what the solution did.} They said \(DA = \frac{325}{22}\), \(DB^2 = \left(\frac{225}{22}\right)^2\). Then:

\[
DP = \frac{(225/22)^2}{325/22}, \quad AP = DP - DA = \frac{2025}{286} - \frac{325}{22} = \frac{-1100}{143}
\]

\textcolor{red}{But length can't be negative, which indicates a mistake in the sign. This suggests that the assumption about the position of P is wrong. Maybe P is between D and A.} So $AP = \frac{1100}{143} = \frac{100}{13} \approx 7.6923$

The solution says \(AP = 22\). That's impossible. The correct value is \(AP = \frac{100}{13}\), hence \(m+n = \boxed{113}\)... \textless/think\textgreater

\#\#\# Decision: \textbf{Continue} \\
\#\#\# Recruit Number: 2 \\
\#\#\# Direction: 1) Verifying the power of point application 2) Re-examining the geometric configuration ... \\
\#\#\# Maximum Tokens: 8192
\end{tcolorbox}

% Problem Solver
\begin{tcolorbox}[
  enhanced,
  colback=white,
  colframe=olivegreen,
  boxrule=0.5pt,
  sharp corners,
  left=1pt, right=1pt, top=1pt, bottom=1pt
]
\textcolor{olivegreen}{\textbf{Problem Solver}} \\
... Power of point \( D \) calculation: $DA \cdot DP = DB^2$

Where: \( DA = \frac{325}{22} \) (distance from \( D \) to \( A \)), \( DB = \frac{225}{22} \) (tangent length from \( D \) to circle). Solving for \( DP \):
\[
DP = \frac{DB^2}{DA} = \frac{\left(\frac{225}{22}\right)^2}{\frac{325}{22}} = \frac{2025}{286}
\]

\red{Critical geometric relationship: Point \( P \) lies between \( D \) and \( A \)}, so: $AP = DA - DP = \frac{325}{22} - \frac{2025}{286} = \frac{100}{13}$

Final result: $AP = \frac{100}{13} \Rightarrow m + n = \boxed{113}$
\end{tcolorbox}
\end{tcolorbox}
\caption{An ``aha" moment in MAS where the CEO agent proactively verifies and corrects the solution provided by the Problem Solver. After identifying an error, the CEO suggests a corrected approach, which the Problem Solver then incorporates into its revised solution.}
\label{fig:aha_moment}
\end{figure}


\subsection{Aha Moment in MAS}
In MAS, we observe that when using M1-32B, agents sometimes exhibit emergent behaviors that actively contribute to validating and refining collaborative processes, even when it is not explicitly required. For example, as illustrated in Figure \ref{fig:aha_moment}, the Problem Solver initially fails to recognize an error in its solution. After reviewing this solution, the CEO agent actively checks its validity and identifies the Problem Solver's error, noting that it results in a negative length for a line segment. The CEO agent then proposes an alternative and correct solution, prompting the Problem Solver to revise its original response accordingly. This collaborative interaction, where one agent assists others by verifying solutions, exploring alternative approaches, and suggesting corrections, occurs even when other agents are unaware of their own mistakes. A plausible reason for this emergent behavior is that the CEO agent, having been trained on multi-agent collaborative reasoning traces and observing other agents' discussions, actively validates and corrects solutions based on learned collaborative patterns and insights gained from the reasoning of other agents.


\subsection{Additional Investigation}
\paragraph{Scaling Collaboration and Reasoning in MAS.} We investigate how scaling collaboration and reasoning affects the performance of M1-32B in MAS by systematically adjusting the total iterations, critic iterations, total agent numbers, and maximum token limits. The results are presented in Figures \ref{fig:ablation_iter} and \ref{fig:scaling_token}. Our observations are as follows: \ding{182} Enhancing collaboration by increasing the interactions among Problem Solvers significantly improves performance. This can be achieved either by raising the critic iteration limit to allow more extensive discussion toward consensus or by increasing the total number of Problem Solvers. However, involving too many Problem Solvers may reduce performance due to divergent discussions among agents. Additionally, merely increasing the total iterations does not improve MAS performance. \ding{183} Enhancing reasoning capabilities by increasing the maximum allowed tokens per agent effectively improves MAS performance. Furthermore, optimal token limits vary by task; for example, 16384 tokens yield optimal results for AIME2024, whereas 8192 tokens are sufficient for GPQA. This finding supports our motivation for using the CEO agent to dynamically manage token allocation based on specific task requirements.

\begin{figure}[t]
    \centering
    \includegraphics[width=0.98\linewidth]{figures/ablation_iter.png}
    \caption{The effect of scale collaboration in AgentVerse using M1-32B by increasing the total iteration, critic iteration, and total agents involved in the MAS.}
    \label{fig:ablation_iter}
\end{figure}



\begin{center}
\begin{minipage}[t]{0.48\textwidth}
\begin{figure}[H]
    \centering
    \includegraphics[width=\linewidth]{figures/scaling_token.png}
    \caption{Effect of scaling reasoning on AgentVerse using M1-32B by controlling the maximum token usage.}
    \label{fig:scaling_token}
\end{figure}
\end{minipage}
\hfill
\begin{minipage}[t]{0.5\textwidth}
\begin{table}[H]
    \centering
    \small
    \renewcommand{\arraystretch}{1.1}
    \begin{tabular}{lcc}
        \toprule
        \textbf{Setting} & \textbf{AIME2024} & \textbf{GPQA} \\
        \midrule
        Qwen2.5 + SA & 26.7 & 49.0   \\
        Qwen2.5 + MAS & 21.1 & 50.2   \\
        Qwen2.5 + MAS w. CEO & 23.3 & 50.5   \\
        \midrule
        M1-32B + SA  & 46.7 & 58.1  \\
        M1-32B + MAS & 60.0 & 61.1 \\
        M1-32B + MAS w. CEO & \textbf{62.2} & \textbf{62.1} \\
        \bottomrule
    \end{tabular}
    \caption{Comparison of Qwen2.5 and M1-32B used as a single agent (SA), within AgentVerse (MAS), and within the MAS w. CEO.}
    \label{tab:single_agent}
\end{table}
\end{minipage}
\end{center}


\paragraph{Performance of M1-32B as a Single Agent.} We further investigate the performance improvement achieved by using M1-32B within MAS compared to its performance as a single agent. The results are summarized in Table \ref{tab:single_agent}. We observe that employing M1-32B in MAS significantly improves performance compared to its single-agent usage. In contrast, using Qwen2.5 within MAS results in smaller improvements over the single-agent setting, further demonstrating the effectiveness of our proposed method in enhancing MAS performance.

